% ==============================================================================
% Physics Research Poster Template — tikzposter
% ==============================================================================
% A competition-quality poster template for physics research conferences.
% Color scheme: Deep Purple & Emerald Green
% Size: A0 Portrait (841 × 1189 mm)
%
% Compile with:  pdflatex physics_poster.tex
%           or:  lualatex physics_poster.tex  (better font support)
% ==============================================================================

\documentclass[
  25pt,                        % Base font scaling
  a0paper,                     % A0 paper size
  portrait,                    % Portrait orientation
  margin=10mm,                 % Outer margins
  innermargin=15mm,            % Padding inside blocks
  blockverticalspace=15mm,     % Vertical space between blocks
  colspace=15mm,               % Space between columns
  subcolspace=8mm              % Space between subcolumns
]{tikzposter}

% =============================================================================
% PACKAGES
% =============================================================================

\usepackage[T1]{fontenc}
\usepackage{helvet}                    % Helvetica sans-serif
\renewcommand{\familydefault}{\sfdefault}
\usepackage{sfmath}                    % Sans-serif math
\usepackage{amsmath,amssymb}
\usepackage{physics}                   % Physics notation: \bra, \ket, \dv, etc.
\usepackage{siunitx}                   % SI units: \SI{3e8}{\metre\per\second}
% Resolve \qty conflict between physics and siunitx
\AtBeginDocument{\RenewCommandCopy\qty\SI}
\usepackage{graphicx}
\usepackage{booktabs}
\usepackage{enumitem}
\usepackage{qrcode}
% \usepackage{hyperref}               % Optional: uncomment for digital-only posters
% \usepackage{microtype}               % Optional: microtypography (uncomment if needed)
\usepackage{setspace}
\usepackage{multicol}

% Tight lists
\setlist{nosep, leftmargin=1.5em}

% Increase line spacing for readability at distance
\setstretch{1.25}

% =============================================================================
% CUSTOM COLOR SCHEME — Deep Purple & Emerald Green
% =============================================================================

\definecolorstyle{PhysicsGreenPurple}{
  % Define base colors
  \definecolor{DeepPurple}{RGB}{75,0,130}
  \definecolor{MediumPurple}{RGB}{120,60,180}
  \definecolor{LightPurple}{RGB}{220,208,240}
  \definecolor{EmeraldGreen}{RGB}{0, 15, 13}
  \definecolor{DarkGreen}{RGB}{0, 15, 13}
  \definecolor{LightGreen}{RGB}{0, 15, 13}
  \definecolor{NearBlack}{RGB}{33,33,33}
  \definecolor{DarkGray}{RGB}{80,80,80}
}{
  % Background Colors
  \colorlet{backgroundcolor}{white}
  \colorlet{framecolor}{DeepPurple}
  % Title Colors
  \colorlet{titlefgcolor}{white}
  \colorlet{titlebgcolor}{DeepPurple}
  % Block Colors
  \colorlet{blocktitlebgcolor}{DeepPurple}
  \colorlet{blocktitlefgcolor}{white}
  \colorlet{blockbodybgcolor}{white}
  \colorlet{blockbodyfgcolor}{NearBlack}
  % Inner Block Colors
  \colorlet{innerblocktitlebgcolor}{EmeraldGreen}
  \colorlet{innerblocktitlefgcolor}{white}
  \colorlet{innerblockbodybgcolor}{LightGreen}
  \colorlet{innerblockbodyfgcolor}{NearBlack}
}

\usecolorstyle{PhysicsGreenPurple}

% Use a clean block style
\useblockstyle{Default}

% Remove the tikzposter attribution logo
\tikzposterlatexaffectionproofoff

% =============================================================================
% CUSTOM HIGHLIGHT BOX (green accent)
% =============================================================================

\usepackage{tcolorbox}
\tcbuselibrary{skins}

\newtcolorbox{highlightbox}[1][]{%
  enhanced,
  colback=LightGreen,
  colframe=EmeraldGreen,
  coltitle=white,
  fonttitle=\bfseries\large,
  boxrule=2pt,
  arc=4mm,
  left=8mm, right=8mm, top=4mm, bottom=4mm,
  #1
}

\newtcolorbox{purplebox}[1][]{%
  enhanced,
  colback=LightPurple,
  colframe=DeepPurple,
  coltitle=white,
  fonttitle=\bfseries\large,
  boxrule=2pt,
  arc=4mm,
  left=8mm, right=8mm, top=4mm, bottom=4mm,
  #1
}

% =============================================================================
% POSTER METADATA — CUSTOMIZE BELOW
% =============================================================================

\title{\textbf{Your Research Title: Concise and Descriptive}}

\author{%
  Author One\textsuperscript{1},
  Author Two\textsuperscript{2},
  \underline{Presenting Author}\textsuperscript{1}%
}

\institute{%
  \textsuperscript{1}Department of Physics, University Name, City, Country \quad
  \textsuperscript{2}Research Institute Name, City, Country%
}

% Logos flanking the title (replace with your institution logos)
% Uncomment and update paths when you have logo files:
% \titlegraphic{%
%   \includegraphics[width=0.08\textwidth]{figures/logo_left.pdf}%
%   \hfill
%   \includegraphics[width=0.08\textwidth]{figures/logo_right.pdf}%
% }

% =============================================================================
% BEGIN DOCUMENT
% =============================================================================

\begin{document}

\maketitle

% =============================================================================
% MAIN CONTENT — 3 COLUMN LAYOUT
% =============================================================================

\begin{columns}

% =============================================================================
% LEFT COLUMN
% =============================================================================
\column{0.33}

\block{Introduction}{
  \textbf{Background}

  Provide 1--2 sentences of context establishing the importance of your
  research area and the physical system under study.

  \vspace{0.5cm}

  \textbf{Motivation}

  What open question or discrepancy does this work address? Why is it
  significant for the broader physics community?

  \vspace{0.5cm}

  \textbf{Objective}

  \begin{highlightbox}
    We investigate \textbf{[physical phenomenon]} using
    \textbf{[technique/method]} to determine \textbf{[key quantity or
    relationship]}.
  \end{highlightbox}
}

\block{Theoretical Framework}{
  \textbf{Key Equations}

  The dynamics are governed by the Hamiltonian:
  \begin{equation}
    \hat{H} = -\frac{\hbar^2}{2m}\nabla^2 + V(\mathbf{r})
    \label{eq:hamiltonian}
  \end{equation}

  The time evolution of the system is given by:
  \begin{equation}
    i\hbar \pdv{\Psi}{t} = \hat{H}\Psi
    \label{eq:schrodinger}
  \end{equation}

  \vspace{0.3cm}

  \textbf{Key Parameters}
  \begin{itemize}
    \item Coupling constant: $g = \qty{1.43e-3}{}$
    \item Temperature range: $T = \qtyrange{4.2}{300}{\kelvin}$
    \item Energy scale: $E_0 = \qty{13.6}{\electronvolt}$
  \end{itemize}
}

\block{Experimental Setup / Methods}{
  \textbf{Apparatus}
  \begin{itemize}
    \item Detector: Description and key specifications
    \item Source: Beam energy, luminosity, or flux
    \item Data acquisition: Sampling rate, channels
  \end{itemize}

  \vspace{0.5cm}

  % Replace with your own methods diagram
  \begin{center}
    \textcolor{DarkGray}{\rule{0.85\linewidth}{0.4pt}}\\[0.5cm]
    {\Large\textcolor{DarkGray}{\textit{[Experimental schematic or
    methods flowchart]}}}\\[0.3cm]
    {\large\texttt{figures/setup\_diagram.pdf}}\\[0.5cm]
    \textcolor{DarkGray}{\rule{0.85\linewidth}{0.4pt}}
  \end{center}

  \vspace{0.3cm}

  \textbf{Analysis}
  \begin{itemize}
    \item Monte Carlo simulations ($N = 10^6$ events)
    \item $\chi^2$ fitting with systematic uncertainties
    \item Software: ROOT, Python (NumPy, SciPy)
  \end{itemize}
}

% =============================================================================
% MIDDLE COLUMN
% =============================================================================
\column{0.33}

\block{Results: Primary Finding}{
  Brief statement of the main result. What is the key observation
  or measurement?

  \vspace{0.5cm}

  % Replace with your primary result figure
  \begin{center}
    \textcolor{DarkGray}{\rule{0.9\linewidth}{0.4pt}}\\[1cm]
    {\Large\textcolor{DarkGray}{\textit{[Primary result figure ---
    spectrum, phase diagram, or data plot]}}}\\[0.5cm]
    {\large\texttt{figures/result\_main.pdf}}\\[1cm]
    \textcolor{DarkGray}{\rule{0.9\linewidth}{0.4pt}}
  \end{center}

  \vspace{0.3cm}

  \textbf{Figure~1:} Measured spectrum showing the resonance at
  $E = \qty{125.3 \pm 0.4}{\giga\electronvolt}$.
  Error bars represent statistical uncertainties;
  shaded band shows systematic uncertainty.
  The fit yields $\chi^2/\text{ndf} = 1.03$.
}

\block{Results: Secondary Finding}{
  Description of the second key result or supporting measurement.

  \vspace{0.5cm}

  % Replace with your secondary result figure
  \begin{center}
    \textcolor{DarkGray}{\rule{0.9\linewidth}{0.4pt}}\\[1cm]
    {\Large\textcolor{DarkGray}{\textit{[Secondary result ---
    comparison, correlation, or trend]}}}\\[0.5cm]
    {\large\texttt{figures/result\_secondary.pdf}}\\[1cm]
    \textcolor{DarkGray}{\rule{0.9\linewidth}{0.4pt}}
  \end{center}

  \vspace{0.3cm}

  \textbf{Figure~2:} Temperature dependence of the order parameter.
  Solid line: theoretical prediction. Data points: experimental
  measurements ($n = 50$ per temperature).
  Agreement within $2\sigma$ across the full range.
}

% =============================================================================
% RIGHT COLUMN
% =============================================================================
\column{0.33}

\block{Results: Additional Analysis}{
  Third key result, systematic study, or validation measurement.

  \vspace{0.5cm}

  % Replace with your third figure
  \begin{center}
    \textcolor{DarkGray}{\rule{0.9\linewidth}{0.4pt}}\\[1cm]
    {\Large\textcolor{DarkGray}{\textit{[Additional analysis ---
    systematic uncertainties, simulation comparison]}}}\\[0.5cm]
    {\large\texttt{figures/result\_analysis.pdf}}\\[1cm]
    \textcolor{DarkGray}{\rule{0.9\linewidth}{0.4pt}}
  \end{center}

  \vspace{0.3cm}

  \textbf{Figure~3:} Comparison of measured cross-section with
  Standard Model prediction (dashed) and BSM scenario (dotted).
  The data favour the SM with $p > 0.95$.

  \vspace{0.5cm}

  \begin{highlightbox}[title=Summary of Measurements]
    \begin{tabular}{@{}lcc@{}}
      \toprule
      \textbf{Observable} & \textbf{Measured} & \textbf{Predicted} \\
      \midrule
      $\sigma$ [fb] & $48.3 \pm 2.1$ & $47.8 \pm 1.5$ \\
      $m$ [\unit{\giga\electronvolt}] & $125.3 \pm 0.4$ & $125.1$ \\
      $\Gamma$ [\unit{\mega\electronvolt}] & $4.1 \pm 0.6$ & $4.07$ \\
      \bottomrule
    \end{tabular}
  \end{highlightbox}
}

\block{Conclusions}{
  \textbf{Key Findings}
  \begin{itemize}
    \item \textbf{Finding 1:} Primary result with quantified significance
    \item \textbf{Finding 2:} Secondary result and its physical implication
    \item \textbf{Finding 3:} Agreement or tension with theoretical predictions
  \end{itemize}

  \vspace{0.5cm}

  \textbf{Outlook}
  \begin{itemize}
    \item Planned improvements (higher statistics, new detector, etc.)
    \item Broader impact on the field
    \item Next experimental campaign or theoretical extension
  \end{itemize}
}

\block{Scan for More}{
  \begin{center}
    \qrcode[height=5cm]{https://doi.org/10.1234/your-paper}\\[0.5cm]
    {\large Full paper, data, and analysis code}
  \end{center}
}

\end{columns}

% =============================================================================
% FOOTER — Full Width
% =============================================================================

\block[titleoffsety=-8mm,bodyoffsety=-8mm]{}{
  \footnotesize
  \begin{minipage}{0.63\textwidth}
    \textbf{References}
    \begin{enumerate}[label={[\arabic*]}]
      \item A.~Author \textit{et al.} (2025). Title of paper.
            \textit{Phys.\ Rev.\ Lett.}\ \textbf{130}, 012345.
      \item B.~Author \textit{et al.} (2024). Title of paper.
            \textit{JHEP}\ \textbf{04}, 067.
      \item C.~Author \textit{et al.} (2024). Title of paper.
            \textit{Nature Physics}\ \textbf{20}, 123--130.
      \item D.~Author \textit{et al.} (2023). Title of paper.
            \textit{Eur.\ Phys.\ J.\ C}\ \textbf{83}, 456.
    \end{enumerate}

    \vspace{0.3cm}

    \textbf{Acknowledgements} \quad
    This work was supported by [Funding Agency] (Grant No.\ XXXX).
    We thank [Collaborators/Facilities] for [contribution].

    \vspace{0.2cm}

    \textbf{Contact:} \quad
    presenter.email@university.edu
    \quad $\vert$ \quad
    labname.university.edu
  \end{minipage}%
  \hfill
  \begin{minipage}{0.28\textwidth}
    \raggedleft
    {\normalsize Conference Name 2026}\\
    Location, Dates\\
    Poster \#XXX\\[0.5cm]
    \textcolor{EmeraldGreen}{\rule{0.8\linewidth}{2pt}}
  \end{minipage}
}

\end{document}
